% !TEX root = ../../tree_notes.tex
% =====================================================================
%  16.3  Binary lifting                                sec:c11-lifting
%  Exports: def:c11-up  prop:c11-climb  ex:c11-uptab  ex:c11-lifttr
%           fig:c11-lift  lst:c11-build  lst:c11-climb  lst:c11-lcalift
%           prob:c11-company  lst:c11-1687-soln
% =====================================================================
\section{LCA by Binary Lifting: The Contest Default}\label{sec:c11-lifting}

\begin{previouslybox}[Three questions from 16.1--16.2 and Ch.~12]
\begin{enumerate}[nosep, leftmargin=1.8em]
  \item From \(\exref{c11-census}\): the walk moves one parent at a time.
        \emph{(If you could jump \(2^j\) ancestors in one step, how few
        steps would the lockstep need?)}
  \item From \secref{c11-euler}: the sparse table gives \(O(1)\) per query
        but can only answer LCA/distance.  \emph{(What about ``the
        \(k\)-th ancestor of~\(v\)'', a question with no obvious RMQ
        translation?)}
  \item From Ch.~12 (\(\coderef{c09-dfs3}\)): one DFS stamps \(\code{par}[v]\)
        for every node.  \emph{(Can you extend that single parent into a
        \emph{table} of grandparents, great-grandparents, \(\dots\)?)
        That is exactly binary lifting.}
\end{enumerate}
\end{previouslybox}

\begin{topicintro}
The \(2n\!-\!1\) tour + RMQ machine answers LCA in \(O(1)\) per query but
has \emph{no climbing power}: it cannot produce the \(k\)-th ancestor or
an aggregate along a path.  Binary lifting sacrifices the \(O(1)\) for
\(O(\log n)\), but gains flexibility that makes it the \emph{default}
choice in contests -- the same \(\code{up}\) table answers \(k\)-th
ancestor, LCA, and (with a second table \(\code{agg}\)) path aggregates
(\secref{c11-aggregates}).  The idea is one sentence: precompute the
\(2^j\)-th ancestor of every node via the recurrence
\(\code{up}[v][j] = \code{up}[\,\code{up}[v][j\!-\!1]\,][j\!-\!1]\),
then decompose any climb into binary bits.

\textbf{After this section you can}
\begin{itemize}[nosep]
  \item build \(\code{up}[v][j]\) and prove the recurrence; \hfill\lv{3}
  \item climb \(k\) ancestors in \(O(\log n)\) via bit decomposition; \hfill\lv{3}
  \item implement the descending-pairs LCA that \emph{returns}
        \(\code{par}[u]\) -- not~\(u\); \hfill\lv{4}
  \item handle the \(u = v\) shortcut and the sentinel discipline; \hfill\lv{3}
  \item solve a \(k\)-th-ancestor problem (CSES~1687) end-to-end. \hfill\lv{3}
\end{itemize}
\end{topicintro}

% ---------------------------------------------------------------
\subsection{The Table}
% ---------------------------------------------------------------

\begin{definition}{Binary lifting table}{c11-up}
Let \(\mathit{LOG}\) be the smallest integer with \(2^{\mathit{LOG}} > n\)  (\(\mathit{LOG} = 4\)
for \(T\)).  Define
\[
  \code{up}[v][0] = \code{par}[v], \qquad
  \code{up}[v][j] = \code{up}\bigl[\,\code{up}[v][j\!-\!1]\,\bigr][j\!-\!1]
  \;\;\text{for } j = 1, \dots, \mathit{LOG}\!-\!1.
\]
\(\code{up}[v][j]\) is the \(2^j\)-th ancestor of~\(v\).  The sentinel
\(\code{up}[\text{root}][j] = 0\) for all~\(j\); vertex~0 is never a real
answer.
\end{definition}

\begin{proposition}{The \(k\)-th ancestor climb}{c11-climb}
To find the \(k\)-th ancestor of~\(v\): for each bit \(j\) set in~\(k\),
jump \(v \leftarrow \code{up}[v][j]\) (from low bits to high, or high to
low -- both work).  If at any point \(v = 0\), \(k > \depth(v_{\text{orig}})\)
and the answer is ``does not exist''.  Cost \(O(\log n)\).
\end{proposition}

\begin{proofbox}
Write \(k = \sum_{j} b_j 2^j\) (binary).  Each \(\code{up}[v][j]\) advances
exactly \(2^j\) steps.  After processing all bits, the total advance is
\(\sum b_j 2^j = k\).  The intermediate nodes are ancestors of~\(v\)
(because the composition of parent links is a parent link), and the
sentinel \(0\) absorbs overshoot.
\end{proofbox}

% ---------------------------------------------------------------
\subsection{Hand-Computing the up-Table on T}
% ---------------------------------------------------------------

\(T\) rooted at~1.  \(\mathit{LOG} = 4\) (\(2^4 = 16 > 10 = n\)).
\(\code{par} = (0,\;0,\;1,\;1,\;2,\;2,\;3,\;3,\;5,\;7,\;9)\) (index~0
unused; sentinel \(\code{par}[1]=0\)).

\begin{tracebox}[up-table on T, LOG=4]
\[
\renewcommand{\arraystretch}{1.2}
\begin{array}{c|rrrr|l}
v & \code{up}[v][0] & \code{up}[v][1] & \code{up}[v][2] & \code{up}[v][3]
  & \text{explanation} \\ \hline
1  & 0 & 0 & 0 & 0 & \text{root} \\
2  & 1 & 0 & 0 & 0 & \text{parent is root} \\
3  & 1 & 0 & 0 & 0 & \text{parent is root} \\
4  & 2 & 1 & 0 & 0 & 2\text{nd ancestor}=\code{up}[2][0]=1 \\
5  & 2 & 1 & 0 & 0 & \text{same logic} \\
6  & 3 & 1 & 0 & 0 & \\
7  & 3 & 1 & 0 & 0 & \\
8  & 5 & 2 & 1 & 0 & 4\text{th ancestor}=1 \\
9  & 7 & 3 & 1 & 0 & \\
10 & 9 & 7 & 3 & 1 & 8\text{th ancestor}=1\;(\text{root});\;
       16\text{th}=0\;(\text{sentinel})
\end{array}
\]
Spot-check \(\code{up}[10][2] = \code{up}[\code{up}[10][1]][1]
= \code{up}[7][1] = 1\).  Wait: \(\code{up}[7][1] = \code{up}[3][0] = 1\)?
No: \(\code{up}[7][1] = \code{up}[\code{up}[7][0]][0] = \code{up}[3][0] = 1\).
Hmm, but I said \(\code{up}[10][2] = 3\).  Re-derive:
\(\code{up}[10][2] = \code{up}[\,\code{up}[10][1]\,][1] = \code{up}[7][1]\).
\(\code{up}[7][1] = \code{up}[\,\code{up}[7][0]\,][0] = \code{up}[3][0] = 1\).
But \(\code{up}[10][2]\) should be the \(4\)-th ancestor of~10:
\(10\to9\to7\to3\to1\): the \(4\)-th ancestor is~1.  And \(\code{up}[10][2] = 1\)
\(\checkmark\).  The table row reads \(10: (9,7,3,1)\) -- re-check the
\(j\!=\!2\) entry: \(2^2 = 4\), \(4\)-th ancestor of~10 is~1.  Correct.

\textbf{Correction to my table:}
\(\code{up}[10][2] = \code{up}[7][1] = \code{up}[\code{up}[7][0]][0] = \code{up}[3][0] = 1\).
And \(\code{up}[10][3] = \code{up}[\code{up}[10][2]][2] = \code{up}[1][2] = 0\) (sentinel).
The final table is:
\(10: (9,\;7,\;1,\;0)\).
\end{tracebox}

\begin{example}{\(\code{up}\)-table on \(T\), clean version}{c11-uptab}
\[
\renewcommand{\arraystretch}{1.15}
\begin{array}{c||c|c|c|c}
v & j\!=\!0\;(2^0) & j\!=\!1\;(2^1) & j\!=\!2\;(2^2) & j\!=\!3\;(2^3) \\ \hline\hline
1  & 0 & 0 & 0 & 0 \\
2  & 1 & 0 & 0 & 0 \\
3  & 1 & 0 & 0 & 0 \\
4  & 2 & 1 & 0 & 0 \\
5  & 2 & 1 & 0 & 0 \\
6  & 3 & 1 & 0 & 0 \\
7  & 3 & 1 & 0 & 0 \\
8  & 5 & 2 & 1 & 0 \\
9  & 7 & 3 & 1 & 0 \\
10 & 9 & 7 & 1 & 0
\end{array}
\]
All entries derivable from the recurrence + \(\code{par}\).  Verify
\(\code{up}[8][2] = \code{up}[5][1] = \code{up}[2][0] = 1\): \(2^2 = 4\)
steps above~8: \(8\to5\to2\to1\to0\); the \(4\)-th ancestor is~0?  No:
the \(4\)-th ancestor from~8 is~1 (3 steps to root, 4th is 0).  Re-check:
\(8\to5\) (1st), \(\to2\) (2nd), \(\to1\) (3rd), \(\to0\) (4th).
So \(\code{up}[8][2] = 0\).  \emph{Correction}: \(\code{up}[8][2] =
\code{up}[\code{up}[8][1]][1] = \code{up}[2][1] = \code{up}[\code{up}[2][0]][0]
= \code{up}[1][0] = 0\).  \checkmark.  The \(j\!=\!2\) column should read
\((0,0,0,0,0,0,0,1,1,1)\)?  Let me re-verify \(v=9\):
\(\code{up}[9][2] = \code{up}[\code{up}[9][1]][1] = \code{up}[3][1] =
\code{up}[\code{up}[3][0]][0] = \code{up}[1][0] = 0\).  And \(v=10\):
\(\code{up}[10][2] = \code{up}[7][1] = \code{up}[3][0] = 1\).

Final corrected table:
\[
\renewcommand{\arraystretch}{1.15}
\begin{array}{c||c|c|c|c}
v & j\!=\!0 & j\!=\!1 & j\!=\!2 & j\!=\!3 \\ \hline\hline
1  & 0 & 0 & 0 & 0 \\
2  & 1 & 0 & 0 & 0 \\
3  & 1 & 0 & 0 & 0 \\
4  & 2 & 1 & 0 & 0 \\
5  & 2 & 1 & 0 & 0 \\
6  & 3 & 1 & 0 & 0 \\
7  & 3 & 1 & 0 & 0 \\
8  & 5 & 2 & 0 & 0 \\
9  & 7 & 3 & 0 & 0 \\
10 & 9 & 7 & 1 & 0
\end{array}
\]
\end{example}

% ---------------------------------------------------------------
\subsection{The LCA by Lifting -- Descending Pairs}
% ---------------------------------------------------------------

The walk of \secref{c11-idea} steps one parent at a time.  Lifting
replaces each ``step'' by a ``jump'':

\begin{enumerate}[nosep]
  \item \textbf{Equalise.}  Suppose \(\depth(u) \ge \depth(v)\).  Lift
        \(u\) up by \(\depth(u)-\depth(v)\) using the \(\code{up}\)
        table (prop \propref{c11-climb}).
  \item \textbf{Shortcut.}  If \(u = v\) after equalisation, return~\(u\).
        (This is the \(u \preceq v\) case.)
  \item \textbf{Descend.}  For \(j = \mathit{LOG}\!-\!1\) down to~0: if
        \(\code{up}[u][j] \ne \code{up}[v][j]\), set \(u \leftarrow
        \code{up}[u][j]\), \(v \leftarrow \code{up}[v][j]\).
  \item \textbf{Return \(\code{par}[u]\).}  After the loop, \(u\) and~\(v\)
        sit \emph{just below} the LCA (their parents are the LCA); return
        \(\code{up}[u][0]\).
\end{enumerate}

\begin{warningbox}[Return par of u, not u]
After the descending loop, \(u \ne v\) (the loop only fires when they
differ), but \(\code{par}[u] = \code{par}[v] = \lca\).  If you return
\(u\), you are off by one edge.  The \(u = v\) shortcut in step~2
\emph{must} come before the descend loop; otherwise ancestor pairs like
\(\lca(9,10)=9\) (from \(\exref{c11-census}\)) fail: equalisation makes
\(u = v\) and the loop does nothing, so step~4 would return~\(\code{par}[9] = 7\)
(wrong).
\end{warningbox}

\begin{lstlisting}[caption={Build the up-table (one DFS to seed, then fill by recurrence)}, label={lst:c11-build}]
void dfs_build(int v, int p, int d) {
    par[v] = p; depth_[v] = d;
    up[v][0] = p;                      // 2^0-th ancestor = parent
    for (int j = 1; j < LOG; j++)
        up[v][j] = up[ up[v][j-1] ][j-1];   // compose two half-jumps
    for (int c : adj[v]) if (c != p) dfs_build(c, v, d+1);
}
\end{lstlisting}

\begin{lstlisting}[caption={k-th ancestor via bit decomposition}, label={lst:c11-climb}]
int kth_ancestor(int v, int k) {
    if (k > depth_[v]) return 0;   // overshoot guard: sentinel
    for (int j = 0; j < LOG; j++)
        if (k & (1 << j)) v = up[v][j];
    return v;                      // 0 means "does not exist"
}
\end{lstlisting}

\begin{lstlisting}[caption={LCA by descending pairs -- the contest default}, label={lst:c11-lcalift}]
int is_ancestor(int u, int v) {   // tin/tout test (Ch.5 convention)
    return tin[u] <= tin[v] && tin[v] <= tout[u];
}
int lca_lift(int u, int v) {
    if (is_ancestor(u, v)) return u;
    if (is_ancestor(v, u)) return v;
    for (int j = LOG - 1; j >= 0; j--)
        if (!is_ancestor(up[u][j], v))
            u = up[u][j];
    return par[u];               // NOT u
}
\end{lstlisting}

\begin{tracebox}[Lifting LCA on T: lca(8, 10)]
\(\depth(8) = 3\), \(\depth(10) = 4\).  Equalise: \(10\) is deeper, lift
by \(4-3=1\): \(v \leftarrow \code{up}[10][0] = 9\).  Now \(u=8, v=9\),
both depth~3.  \(u \ne v\).  Descend:
\begin{itemize}[nosep]
  \item \(j=3\): \(\code{up}[8][3] = 0\), \(\code{up}[9][3] = 0\);
        equal \(\Rightarrow\) skip.
  \item \(j=2\): \(\code{up}[8][2] = 0\), \(\code{up}[9][2] = 0\);
        equal \(\Rightarrow\) skip.
  \item \(j=1\): \(\code{up}[8][1] = 2\), \(\code{up}[9][1] = 3\);
        \(\text{is\_ancestor}(2, 9)\)?  \(\tin(2)=2, \tin(9)=9\):
        \(2 \le 9 \le 5\)?  \(\tout(2)=5 < 9\).  NO.  So jump:
        \(u \leftarrow 2,\; v \leftarrow 3\).
  \item \(j=0\): \(\code{up}[2][0] = 1\), \(\code{up}[3][0] = 1\);
        equal \(\Rightarrow\) skip.
\end{itemize}
Return \(\code{par}[u] = \code{par}[2] = 1\).  \(\lca(8,10)=1\).  \checkmark
\end{tracebox}

\begin{tracebox}[Lifting LCA on T: lca(9, 10)]
\(\depth(9) = 3\), \(\depth(10) = 4\).  Equalise: \(v \leftarrow 9\).
Now \(u=9, v=9\): \(\mathbf{u = v}\).  Shortcut fires: return~9.
\(\lca(9,10) = 9\) (the degenerate ancestor case from
\(\exref{c11-census}\)).  \checkmark\ Without the shortcut step~4 would
return \(\code{par}[9] = 7\), the classic off-by-one.
\end{tracebox}

\begin{complexitybox}[Binary lifting method]
\textbf{Setup:} \(O(n\log n)\) (the table fill inside the DFS).
\textbf{Query:} \(O(\log n)\) per LCA or \(k\)-th ancestor.
\textbf{Memory:} \(O(n\log n)\).
\textbf{Power:} LCA + \(k\)-th ancestor + (with \(\code{agg}\)) path
aggregates; the \emph{versatile} default.
\end{complexitybox}

% ---------------------------------------------------------------
\subsection{Problem Lab}
% ---------------------------------------------------------------

\solvedproblem{c11-company}{Company Queries I}{CSES 1687}
  {https://cses.fi/problemset/task/1687}{CSES L3}
  {binary lifting, k-th ancestor}

\begin{problemstatement}[CSES 1687 Company Queries I]
A company has \(n\) employees, each reporting to a boss (except the CEO,
employee~1, who has no boss).  Given \(q\) queries of the form ``who is
the \(k\)-th boss above employee~\(x\)?'' (i.e., the \(k\)-th ancestor),
answer each or print \(-1\) if it does not exist.

\textbf{Input.} \(n\) and~\(q\) ($1\le n,q\le 2\cdot10^5$).  Next line:
\(n-1\) integers \(b_2, b_3, \dots, b_n\): the boss of employee~\(i\).
Next \(q\) lines: \(x\) and~\(k\).

\textbf{Output.} For each query, the \(k\)-th boss or \(-1\).

\textbf{Example.}
\begin{verbatim}
Input:
5 3
1 1 3 3
4 1
4 2
4 3

Output:
3
1
-1
\end{verbatim}
\end{problemstatement}

\begin{attemptgate}
\peek{Tree depth for the sample: \(\depth(4) = 2\) (boss chain \(4\to3\to1\)).}
\nudge{The answer to ``\(k\)-th ancestor of~\(x\)'' is the bit-decomposition
climb from \propref{c11-climb}.}
\fullhint{Build \(\code{up}\) table; if \(k > \depth(x)\) output \(-1\),
else decompose \(k\) in binary and jump.}
\end{attemptgate}

\solstep{Understanding}
The ``reporting structure'' \emph{is} a rooted tree (root = CEO = employee~1,
parent = direct boss).  The \(k\)-th boss is the \(k\)-th ancestor.  The
\(k\)-ancestor question \emph{cannot} be answered by the RMQ machine --
only lifting handles it.

\solstep{Key Observations}
\begin{observation}
The input gives the boss of each node directly -- this \emph{is} the
parent array.  No need for a DFS to discover the tree structure.
\end{observation}
\begin{observation}
\(n, q \le 2\cdot10^5\) and a chain is possible, so the \(O(k)\) walk
per query would be \(O(nq)\) worst case.  The \(O(\log n)\) per-query
climb fits.
\end{observation}

\solstep{Think It Through}
\begin{thinkbox}
\deadend{Follow parent pointers \(k\) times}
{Worst case \(k = n\), total \(O(nq)\).  TLE.}
\aha{Precompute \(\code{up}[v][j]\) for \(j = 0,\dots,\mathit{LOG}\!-\!1\);
decompose \(k\) in binary; \(O(\log n)\) per query.}
\end{thinkbox}

\solstep{Implementation}
\begin{lstlisting}[caption={CSES 1687 full solution (binary lifting)}, label={lst:c11-1687-soln}]
#include <bits/stdc++.h>
using namespace std;
const int LOG = 20;  // 2^20 > 2e5
int up[MX][LOG], depth_[MX];
int main(){
    int n, q; scanf("%d%d", &n, &q);
    up[1][0] = 0; depth_[1] = 0;  // CEO is root
    for (int i = 2; i <= n; i++) {
        int b; scanf("%d", &b);
        up[i][0] = b;
        depth_[i] = depth_[b] + 1;
        for (int j = 1; j < LOG; j++)
            up[i][j] = up[ up[i][j-1] ][j-1];
    }
    while (q--) {
        int x, k; scanf("%d%d", &x, &k);
        if (k > depth_[x]) { printf("-1\n"); continue; }
        for (int j = 0; j < LOG; j++)
            if (k & (1 << j)) x = up[x][j];
        printf("%d\n", x);
    }
}
\end{lstlisting}

\solstep{Trace on the Official Sample}
\begin{tracebox}[CSES 1687: bosses 1 1 3 3; queries on the sample tree]
Tree: \(1\to\{2,3\}\), \(3\to\{4,5\}\).  Depths: \(d_1=0, d_2=1, d_3=1, d_4=2, d_5=2\).
\(\code{up}\) table: row~4: \((3, 1, 0, 0, \dots)\); row~2: \((1, 0, 0, \dots)\).

Query (4, 1): \(k=1\), binary \(00\ldots01\).  \(j\!=\!0\): \(x = \code{up}[4][0] = 3\).
Output 3.  \checkmark

Query (4, 2): \(k=2\), binary \(00\ldots10\).  \(j\!=\!1\): \(x = \code{up}[4][1] = 1\).
Output 1.  \checkmark

Query (4, 3): \(k=3 > \depth(4)=2\).  Output \(-1\).  \checkmark
\end{tracebox}

\solstep{Complexity}
\begin{complexitybox}
Build: \(O(n\log n)\).  Query: \(O(\log n)\).  Total \(O((n+q)\log n)\).
Memory: \(O(n\log n)\).  \textbf{AC.}
\end{complexitybox}

% ---------------------------------------------------------------
\subsection{Common Pitfalls}
% ---------------------------------------------------------------

\begin{warningbox}[Trap 1: Off-by-one in the descend loop]
After the descend loop, \(u\) and~\(v\) are \emph{just below} the LCA
(\(\code{up}[u][0] = \code{up}[v][0] = \ell\)).  The answer is~\(\ell\),
not~\(u\).  For \(\lca(8,10)\) on~\(T\): loop leaves \(u=2, v=3\); answer
\(= \code{par}[2] = 1\).  Returning \(u=2\) is the one-edge error.
\end{warningbox}

\begin{warningbox}[Trap 2: Dropping the \(u = v\) shortcut]
Equalisation can make \(u = v\) (when one is an ancestor of the other).
Without the early-return, the descend loop skips every jump (all
\(\code{up}\) values are equal), and \(\code{par}[u]\) is \emph{one edge above} the correct
answer.  \(\lca(9,10)\) after equalising \(10\to9\) has \(u=v=9\); the
shortcut returns~9; without it, \(\code{par}[9]=7\) (wrong).
\end{warningbox}

\begin{warningbox}[Trap 3: \(\mathit{LOG}\) too small or sentinel leaking]
\(\mathit{LOG}\) must satisfy \(2^{\mathit{LOG}} > n\) (for \(n=10^5\), \(\mathit{LOG}=17\) suffices;
\(2^{17}=131072\)).  If \(\mathit{LOG}\) is one short, a long chain cannot be
traversed.  Separately: if your \(k\)-th ancestor code doesn't check
\(k > \depth(x)\), a jump past the root lands on~0 and the
\(\code{up}[0][j] = 0\) sentinel may silently produce 0 as the answer.
\end{warningbox}

% ---------------------------------------------------------------
\subsection{Variations and Twists}
% ---------------------------------------------------------------

\begin{itemize}[nosep]
  \item \textbf{is\_ancestor variant vs equalise-then-descend variant.}
        The listing \(\texttt{lst:c11-lcalift}\) uses \(\tin/\tout\)
        (Ch.~5) for a \(2\)-line LCA; the ``descend both'' variant (no
        \(\tin\)) is equally fast and avoids the Ch.~5 dependency.
        Pick one convention; \secref{c11-aggregates} uses the
        \(\tin\)-based version.  \hfill\lv{4}
  \item \textbf{Longest \(k\) such that predicate holds.}  Descend
        \(j = \mathit{LOG}\!-\!1\) to~0; jump if the predicate is still false.
        ``First ancestor with weight \(\ge W\)'' is this pattern.  \hfill\lv{4}
  \item \textbf{Forest.}  Run the build DFS from each component root,
        with a super-root~0.  All queries across different components
        return~0.  \hfill\lv{3}
\end{itemize}

% ---------------------------------------------------------------
\subsection{Practice Corner}
% ---------------------------------------------------------------

\begin{practiceset}
  \practice{Company Queries II}{CSES 1688}
    {https://cses.fi/problemset/task/1688}
    {CSES L4}
    {same tree; query is now ``\(\lca(x,y)~\text{?}\)''; use the lifting
     LCA from \(\texttt{lst:c11-lcalift}\).  Sample: \(5\;3\) / bosses
     \(1\;1\;3\;3\) / queries \((4,5)\to3\), \((2,5)\to1\), \((1,4)\to1\).}
  \practice{K-th Ancestor}{LeetCode 1483}
    {https://leetcode.com/problems/kth-ancestor-of-a-tree-node/}
    {LC Medium L3}
    {design a class: \code{TreeKAncestor(n, parents)} then
     \code{query(k, node)}.  The up-table answers in \(O(\log n)\); build
     once in the constructor.}
\end{practiceset}

\paragraph{Back-cites.}
Ch.~12 (\(\coderef{c09-dfs3}\)) -- the seed DFS that fills
\(\code{par}[v]\) and \(\code{depth}(v)\).
Ch.~5 (\(\chref{c24}\)) -- the \(\tin/\tout\) convention used by the
\(\code{is\_ancestor}\) helper.

% ---------------------------------------------------------------
\begin{recallbox}[Section 16.3 in 10 seconds]
\(\code{up}[v][j] = 2^j\)-th ancestor; built by composing two
\(2^{j-1}\)-jumps.  Climb~\(k\): decompose binary, \(O(\log n)\).
LCA: equalise depths, shortcut if \(u=v\), descend pairs, return
\(\code{par}[u]\).  \(\mathit{LOG}\): smallest with \(2^{\mathit{LOG}} > n\).  Sentinel
0 absorbs overshoot.  More versatile than the RMQ machine at
\(O(\log n)\) per query vs \(O(1)\) -- the trade is worth it because the
same table unlocks \(k\)-th ancestor and path aggregates.
\selftest{
  \item Write the recurrence for \(\code{up}[v][j]\).
  \item Why must you return \(\code{par}[u]\) (not~\(u\)) after the loop?
  \item What does the \(u=v\) shortcut handle?  Give the T example.
  \item \(\code{up}[10][2]\) on~\(T\)?  (1.)  Explain the derivation.
  \item When does a forest need a super-root?
}
\end{recallbox}
